We introduced \textbf{RetailBench}, a data-grounded benchmark for evaluating long-horizon decision-making by tool-using LLM agents in realistic retail environments. By integrating interdependent business decisions, dynamic feedback, and delayed consequences, RetailBench captures key complexities of real-world retail operations within a controlled evaluation setting.

Across seven LLMs and 19 model--scaffold configurations, our results reveal substantial challenges in sustaining effective decision-making over extended horizons. Only two rollouts complete the 180-day horizon, and even the best-performing agent substantially underperforms the OR-Composite and Privileged Heuristic references. Our four-stage trajectory analysis further reveals a gap between information acquisition and decision quality: agents often gather relevant evidence but exhibit \textbf{shallow evidence-to-decision integration} and \textbf{limited responsiveness to salient operational anomalies}.
These findings establish RetailBench as a testbed for evaluating and diagnosing long-horizon agent capabilities in realistic economic environments, highlighting the challenges of maintaining coherent and adaptive decision-making under complex operational constraints.